\chapter{Conclusão}
\label{cap:conclusao}

Este trabalho estendeu o MellowDB, um \textit{middleware} de evolução semântica originalmente restrito a um nó único do MongoDB, a ambientes distribuídos, e caracterizou o custo dessa transição. A extensão a \textit{replica sets} mantém a cadeia de versões íntegra e autoritativa no nó primário, escrita sob transação multidocumento, e distribui as leituras de registros aos secundários sem abrir mão dessa cadeia única. A extensão a \textit{sharded clusters} declara a chave de \textit{shard}, materializa as versões sobre uma coleção particionada e trata os documentos que não possuem o campo da chave. Em torno dessas extensões, construiu-se um ferramental reproduzível e uma campanha de experimentos sobre um corpus sintético e um corpus real, executada em um \textit{cluster} na nuvem.

A primeira pergunta indagava se a distribuição altera qual das duas estratégias de evolução predomina. A resposta é que o compromisso entre o \textit{rewrite} e o \textit{preprocess} não tem um vencedor fixo. O \textit{preprocess} lidera em todas as configurações replicadas e particionadas avaliadas, mas a leitura dos secundários encurta a distância, e o \textit{rewrite} aproveita melhor cada nó acrescentado, por ter uma leitura mais pesada que usufrui do paralelismo. A vantagem, de fato, só se inverte quando a cadeia de evolução se aprofunda, caso em que o \textit{rewrite} lê mais rápido que o \textit{preprocess}, pois o custo por consulta do pré-processamento cresce com a profundidade sobre uma varredura sem índice. No corpus real, por outro lado, as consultas amplas sobre milhões de registros tornam o \textit{rewrite} impraticável. A distribuição, portanto, estabelece novos parâmetros a considerar para escolher a estratégia mais adequada.

A segunda pergunta tratava de como a evolução semântica interage com o particionamento físico dos dados. A campanha mostrou que, quando a chave de \textit{shard} é o campo que evolui, aplicar uma operação realoca fisicamente os registros afetados entre os \textit{shards}. O desequilíbrio resultante pode ser irreparável, pois, sob baixa cardinalidade, o valor dominante forma um \textit{jumbo chunk} que o balanceador não consegue dividir nem migrar. No corpus real, o reequilíbrio foi possível, mas custou a movimentação de dezenas de megabytes à medida que a coleção crescia na transição. Chaves alheias ao campo que evolui não deslocam dados, e chaves bem distribuídas, como as \textit{hashed}, paralelizam tanto a operação quanto as consultas amplas. Assim, fica claro que uma chave de \textit{shard} que recai sobre o campo em que a evolução semântica atua pode trazer desvantagens de desempenho.

Entretanto, esses resultados estão delimitados pelas condições dos experimentos. O tamanho de \textit{chunk} foi reduzido em relação ao padrão para manter o balanceador ativo nos volumes de corpus que cabiam no orçamento de tempo da campanha, uma condição declarada que afeta os valores absolutos de reequilíbrio. A avaliação mediu custo e desempenho, e não a consistência em si, um fator importante introduzido na Fundamentação e que ficou fora do escopo. A escala do particionamento limitou-se a sete \textit{shards}, de modo que um eventual cruzamento das estratégias além dessa faixa não foi alcançado, e cada \textit{shard} foi composto por um único nó, o que deixa de fora os efeitos da replicação interna a um \textit{shard}.

Como trabalhos futuros, propõe-se medir a consistência deixada de fora desta avaliação, a defasagem entre réplicas e o comportamento sob falha, tratar os objetos aninhados e a evolução de esquema, e avaliar um sistema de maior escala que combine os dois tipos de distribuição ao mesmo tempo.
